% Part: proof-theory
% Chapter: sequent-calculus
% Section: admissible-derivable

\documentclass[../../../include/open-logic-section]{subfiles}

\begin{document}

\olfileid[tr]{pt}{seq}{adm}

\olsection{Kabul Edilebilir ve Türetilebilir Kurallar}

İspat kuramındaki birçok sonuç, belirli bir kural kullanılarak !!{prove}{}nan
her şeyin o kural olmadan da !!{prove}{}nıp !!{prove}{}namayacağına ilişkin
sonuçlarla ilgilidir ya da bunları kullanır. En iyi bilinen örnek kesme
giderme teoremidir. Bu teorem, \Cut{} içeren bir sekant sisteminin
kurallarıyla !!{prove}{}nabilen her sekantın \Cut{} olmadan da
!!{prove}{}nabildiğini gösterir. Bu özellik tipi çok önemli olduğundan özel
bir adı vardır.

\begin{defn}
Bir~$R$ kuralı, bu kuralla birlikte bir sekant sisteminde ispatlanabilen her
sekant $R$ olmadan da ispatlanabiliyorsa o sistemde \emph{kabul edilebilirdir}.
\end{defn}

\begin{ex}\ollabel{ex:land-deriv}
$\LeftR{\land}$ kuralı \Log{G1c} ve \Log{G3c}'de farklı biçimler alır.
\Log{G1c}'de kuralın iki sürümü vardır: birinin öncülünde ana
!!{formula}~$!A \land !B$'nin sol bağlaşığı~$!A$, ötekinin öncülünde sağ
bağlaşığı~$!B$ bulunur:
\[
\Axiom$ !A, \Gamma \fCenter \Delta$
\RightLabel{$\LeftR{\land}_1$}
\UnaryInf$ !A \land !B, \Gamma \fCenter \Delta$
\DisplayProof
\qquad
\Axiom$!B, \Gamma \fCenter \Delta$
\RightLabel{$\LeftR{\land}_2$}
\UnaryInf$!A \land !B, \Gamma \fCenter \Delta$
\DisplayProof
\]
Buna karşılık \Log{G3c}'nin yalnızca bir kuralı vardır:
\[
\Axiom$ !A, !B, \Gamma \fCenter \Delta$
\RightLabel{$\LeftR{\land}_3$}
\UnaryInf$ !A \land !B, \Gamma \fCenter \Delta$
\DisplayProof
\]
Şimdi ``\Log{G3c}'nin $\LeftR{\land}_3$ kuralı \Log{G1c}'de kabul
edilebilir mi?'' diye sorabiliriz. Evet: $\LeftR{\land}_3$ kullanan her
çıkarımı yalnız \Log{G1c}'nin kurallarıyla doğrudan simüle edebiliriz.
$\LeftR{\land}_1$ ve $\LeftR{\land}_2$'ye ek olarak yapısal
$\LeftR{\Contraction}$ kuralı da gerekir:
\[
\Axiom$!A, !B, \Gamma \fCenter \Delta$
\RightLabel{$\LeftR{\land}_1$}
\UnaryInf$!A \land !B, !B, \Gamma \fCenter \Delta$
\RightLabel{$\LeftR{\land}_2$}
\UnaryInf$!A \land !B, !A \land !B, \Gamma \fCenter \Delta$
\RightLabel{\LeftR{\Contraction}}
\UnaryInf$!A \land !B, \Gamma \fCenter \Delta$
\DisplayProof
\]
\end{ex}

Bir kuralı kullanan çıkarımları başka kurallarla simüle etmek, o kuralın kabul
edilebilir olduğunu göstermenin bir yoludur. Ancak tek yol bu değildir;
simüle edilemediği hâlde kabul edilebilir olan kurallar da vardır.
Simüle edilebilen kurallara \emph{türetilebilir} denir.

\begin{defn}
Bir~$R$ kuralına
\[
\AxiomC{$S_1$}
\AxiomC{$\dots$}
\AxiomC{$S_n$}
\RightLabel{$R$}
\TrinaryInfC{$S$}
\DisplayProof
\]
bir sistemde \emph{türetilebilir} denmesi için, sistemin kuralları ve
aksiyomlarıyla birlikte $S_1$, \dots,~$S_n$ biçimindeki ek başlangıç
sekantlarını kullanarak $S$ sekantının şematik bir !!a{proof}{}ı bulunmalıdır.
\end{defn}

\olref{ex:land-deriv} içindeki !!{proof}, $\LeftR{\land}_3$'ün
\Log{G1c}'de türetilebilir bir kural olduğunu gösterir; çünkü bu !!{proof},
$\LeftR{\land}_3$'ün öncüllerini başlangıç sekantları sayarak, yalnız
\Log{G1c}'nin kurallarıyla kuralın sonucunun şematik bir !!a{proof}{}ıdır.
(\Log{G1c}'nin hiçbir aksiyomunu kullanmamıştır.)

\begin{prop}\ollabel{prop:lor-G3-adm}
\Log{G3c}'nin $\LeftR{\land}$ ve $\RightR{\lor}$ kuralları \Log{G1c}'de
türetilebilirdir.
\end{prop}

\begin{proof}
  $\LeftR{\land}$ ispatı \olref{ex:land-deriv} içinde verilmiştir.
  $\RightR{\lor}$ ispatı alıştırma olarak bırakılmıştır.
\end{proof}

\begin{prob}
\Log{G3c}'nin $\RightR{\lor}$ kuralının \Log{G1c}'de türetilebilir
olduğunu gösterin.
\end{prob}

\begin{prob}
$\liff$'in tanımlı bir !!{operator} olduğunu varsayın: $!A \liff !B$,
$(!A \lif !B) \land (!B \lif !A)$'nın kısaltmasıdır. Şu kuralların
\begin{enumerate}
\item \Axiom$\Gamma, !A, !B \fCenter \Delta$
\Axiom$\Gamma \fCenter \Delta, !A, !B$
\RightLabel{\LeftR{\liff}}
\BinaryInf$!A \liff !B, \Gamma \fCenter \Delta$
\DisplayProof
\item
\Axiom$!A, \Gamma \fCenter \Delta, !B$
\Axiom$!B, \Gamma \fCenter \Delta, !A$
\RightLabel{\RightR{\liff}}
\BinaryInf$\Gamma \fCenter \Delta, !A \liff !B$
\DisplayProof
\end{enumerate}
(a)~\Log{G1c} ve (b)~\Log{G3c}'de türetilebilir olduğunu gösterin.
\end{prob}


\begin{ex}\ollabel{ex:weak-adm}
\Log{G1c}'nin $\RightR{\Weakening}$ kuralı
\[
\Axiom$ \Gamma \fCenter \Delta$
\RightLabel{\RightR{\Weakening}}
\UnaryInf$ \Gamma \fCenter \Delta, !A$
\DisplayProof
\]
\Log{G3c}'de kabul edilebilirdir. İspatın düşüncesi basittir: \Log{G3c}'de
$\Gamma \Sequent \Delta$ öncülünün bir !!a{proof}{}ı~$\pi$ bulunduğunu
varsayın. $\pi$ içindeki her sekantın artbileşenine !!{formula}~$!A$'yı
ekleyin. Aksiyomlar aksiyom olarak, doğru çıkarımlar da doğru çıkarım olarak
kalır. Yeni !!{proof}{}ın son sekantı $\Gamma \Sequent \Delta, !A$'dır.

Ancak $\RightR{\Weakening}$ kuralı \Log{G3c}'de türetilebilir değildir.
Türetilebilir olduğunu, yani \Log{G3c}'nin aksiyom ve kurallarına olası ek
başlangıç sekantı~$\Gamma \Sequent \Delta$ da katılarak $\Gamma \Sequent
\Delta, !A$'nın bir !!a{proof}{}ının bulunabileceğini varsayalım.
$\RightR{\Weakening}$'in türetilebilir olması için gerektiği gibi bu
!!{proof} bütünüyle şematikse $\Gamma$ ile~$\Delta$ silinerek
$\Sequent !A$'nın bir !!a{proof}{}ına dönüştürülebilir. Bu işlem ek başlangıç
sekantı~$\Gamma \Sequent \Delta$'yı boş sekant~$\quad \Sequent \quad$'a
dönüştürür. Boş sekant hiç !!{formula} içermediği, \Log{G3c}'nin bütün
kurallarıysa öncüllerinde en az bir etkin !!{formula} gerektirdiği için boş
sekant bu şematik !!{proof} içindeki hiçbir çıkarımın öncülü olamaz.
Dolayısıyla !!{proof}, ancak ek $\Gamma \Sequent \Delta$ sekantını gerçekten
başlangıç sekantı olarak kullanmıyorsa şematik olabilir. Başka bir deyişle,
yalnız \Log{G3c}'nin aksiyom ve kurallarıyla $\Gamma \Sequent \Delta,
!A$'nın şematik bir türetimi olur ve bu biçimdeki her sekantın \Log{G3c}'de
bir !!a{proof}{}ı bulunduğu gösterilmiş olurdu. Fakat bu olanaksızdır; çünkü
\Log{G3c} sağlamdır ve yalnız her !!{structure} içinde doğru olan sekantları
!!{prove}{}yabilir. $\Gamma = \Delta = \emptyset$ ve $!A \ident \lfalse$
alırsak \Log{G3c}'nin örneğin $\quad \Sequent \lfalse$ sekantını
!!{prove}{}ması gerekirdi; oysa bunu yapamaz.
\end{ex}

Şimdi zayıflatmanın \Log{G3c}'de kabul edilebilir bir kural olduğuna ilişkin
güzel bir tümevarımsal ispat verelim. Aslında sezgisel savın gösterdiği gibi
biraz daha güçlü bir sonuç geçerlidir: $\pi$'nin her sekantının sağ yanına
$!A$ eklemek !!{proof}{}ın yapısını değiştirmez; özellikle yüksekliği
artırmaz. Bu önemli olduğu için ayrı bir tanımı hak eder.

\begin{defn}
Bir~$R$ kuralı
\[
\AxiomC{$S_1$}
\AxiomC{$\dots$}
\AxiomC{$S_n$}
\RightLabel{$R$}
\TrinaryInfC{$S$}
\DisplayProof
\]
bir sistemde, $i = 1$, \dots,~$n$ için $\Proves[h_i] S_i$ olduğunda
$m = \max(h_1, \dots, h_n)$ ile $\Proves[m] S$ oluyorsa
\emph{!!{height} değerini koruyarak kabul edilebilir} (ya da
\emph{!!{hp}-kabul edilebilir}) diye adlandırılır.
\end{defn}

Bir kuralın !!{height} değerini koruyarak kabul edilebilir olması için şu
koşulun yeterli olduğuna dikkat edin: öncüller~$S_i$, !!{height} değeri
$h_i = \pheight{\pi_i}$ olan !!{proof}s~$\pi_i$'ye sahip olduğunda sonucun,
!!{height} değeri~$\pheight{\pi}$, $h_i$'lerin en büyüğünden fazla olmayan bir
!!a{proof}{}ı~$\pi$ bulunsun. Özellikle tek öncül~$S_1$ varsa
$\pheight{\pi} \le \pheight{\pi_1}$ olması yeterlidir.
$\RightR{\Weakening}$ ve $\LeftR{\Weakening}$ durumunda gerçekte
$\pheight{\pi} = \pheight{\pi_1}$ olur; ancak bu gerekli değildir.

\begin{prop}\ollabel{prop:weak-G3c-adm}
\Log{G1c}'nin $\RightR{\Weakening}$ ve $\LeftR{\Weakening}$ kuralları
\Log{G3c}'de !!{height} değerini koruyarak kabul edilebilirdir.
\end{prop}

\begin{proof}
  $\Log{G3c} \Proves \Gamma \Sequent \Delta$ yargısı,
  $\pheight{\pi} = n$ olan bir !!a{proof}~$\pi$ ile geçerliyse,
  $\Gamma \Sequent \Delta, !A$'nın $\pheight{\pi'} = n$ olan
  bir \Log{G3c}-!!{proof}{}ı~$\pi'$ bulunduğunu göstermeliyiz. Bunu $n$
  üzerinde tümevarımla ispatlıyoruz. $n=0$ ise $\pi$ yalnız $\Gamma
  \Sequent \Delta$ aksiyomudur (yani atomik bir $!C$, hem $\Gamma$ hem de
  $\Delta$ içinde geçer); $\Gamma \Sequent \Delta, !A$ da aksiyomdur.
  $n = \pheight{\pi} >0$ ise !!{proof}~$\pi$'nin bir son çıkarımı vardır.
  Kullanılan kurala göre durumları ayırırız. Yalnız bir örnek verelim: son
  çıkarımın $\RightR{\land}$ olduğunu varsayın. O zaman $\Delta = \Delta',
  !B \land !C$ ve son çıkarımın doğrudan alt-!!{proof}s{}ı $\pi_1$ ile
  $\pi_2$ sırasıyla $\Gamma \Sequent \Delta', !B$ ve $\Gamma \Sequent
  \Delta', !C$ son sekantlarına sahiptir. Başka bir deyişle $\pi$ şöyle
  görünür:
  \[
  \AxiomC{}
  \RightLabel{$\pi_1$}
  \Deduce$\Gamma \fCenter \Delta', !B$
  \AxiomC{}
  \RightLabel{$\pi_2$}
  \Deduce$\Gamma \fCenter \Delta', !C$
  \RightLabel{\RightR{\land}}
  \BinaryInf$\Gamma \fCenter \Delta', !B \land !C$
  \DisplayProof
  \]
  Ayrıca $n = \max(\pheight{\pi_1}, \pheight{\pi_2}) + 1$'dir. Dolayısıyla
  $\pheight{\pi_1} < n$ ve $\pheight{\pi_2} < n$ olur ve tümevarım
  varsayımı uygulanır: $\Gamma \Sequent \Delta', !B, !A$ ve $\Gamma
  \Sequent \Delta', !C, !A$'nın !!{proof}s{}ı $\pi_1'$ ile~$\pi_2'$ vardır.
  $\RightR{\land}$ kuralını kullanarak $\Gamma \Sequent \Delta', !B
  \land !C, !A$'nın bir !!a{proof}{}ı~$\pi'$yi elde ederiz:
  \[
  \AxiomC{}
  \RightLabel{$\pi_1'$}
  \Deduce$\Gamma \fCenter \Delta', !B, !A$
  \AxiomC{}
  \RightLabel{$\pi_2'$}
  \Deduce$\Gamma \fCenter \Delta', !C, !A$
  \RightLabel{\RightR{\land}}
  \BinaryInf$\Gamma \fCenter \Delta', !B \land !C, !A$
  \DisplayProof
  \]
  $\pheight{\pi'} = \max(\pheight{\pi_1'}, \pheight{\pi_2'}) + 1 =
  \max(\pheight{\pi_1}, \pheight{\pi_2}) + 1 = \pheight{\pi}$ olur.
  Sol zayıflatma durumu da simetriktir: aynı tümevarımda her sekantın
  önbileşenine $!A$ eklenir ve yükseklik değeri değişmez.
\end{proof}

\olref{prop:weak-G3c-adm} son derece kullanışlıdır. Onun yinelenen
uygulanışı için bir gösterim tanıtacağız: $\pi$, $\Gamma \Sequent \Delta$'nın
bir !!a{proof}{}ı ise $\pi[\Pi \Sequent \Lambda]$, solda $\Pi$ içindeki bütün
!!{formula}s ve sağda $\Lambda$ içindeki bütün !!{formula}s ile
zayıflatmanın kabul edilebilirliğinin uygulanması sonucudur. Bu, $\Gamma,
\Pi \Sequent \Delta, \Lambda$'nın bir !!a{proof}{}ıdır.

\end{document}
